\chapter{実践パターンマッチ指向プログラミング}\label{practice}

\section{リストプログラミング}

\section{ポーカーの役判定}

多重集合に対する非線形パターンを記述できるEgisonを使うとポーカーの役判定が簡潔に記述できる．
それぞれのポーカーの役は直接1つのパターンとして記述できる．
著者がEgisonのマッチ式を設計をしたとき，ポーカーの約判定をするプログラムを簡潔に記述できることを必要条件として設計した．

\begin{lstlisting}[language=egison]
poker cs :=
  match cs as multiset card with
  | card $s $n :: card #s #(n-1) :: card #s #(n-2) :: card #s #(n-3) :: card #s #(n-4) :: _
    -> "Straight flush"
  | card _ $n :: card _ #n :: card _ #n :: card _ #n :: _ :: []
    -> "Four of a kind"
  | card _ $m :: card _ #m :: card _ #m :: card _ $n :: card _ #n :: []
    -> "Full house"
  | card $s _ :: card #s _ :: card #s _ :: card #s _ :: card #s _ :: []
    -> "Flush"
  | card _ $n :: card _ #(n-1) :: card _ #(n-2) :: card _ #(n-3) :: card _ #(n-4) :: []
    -> "Straight"
  | card _ $n :: card _ #n :: card _ #n :: _ :: _ :: []
    -> "Three of a kind"
  | card _ $m :: card _ #m :: card _ $n :: card _ #n :: _ :: []
    -> "Two pair"
  | card _ $n :: card _ #n :: _ :: _ :: _ :: []
    -> "One pair"
  | _ :: _ :: _ :: _ :: _ :: [] -> "Nothing"
\end{lstlisting}

\begin{lstlisting}[language=egison]
suit := algebraicDataMatcher
  | spade
  | heart
  | club
  | diamond

card := algebraicDataMatcher
  | card suit (mod 13)
\end{lstlisting}

\section{SATソルバー}

\section{ゲーム木のパターンマッチ}

